\subsection{Generování cílového kódu}
\label{sec:codegen}

\subsubsection{Backend pro Rust}

Generování cílového kódu pro Rust je implementováno dvěma cestami.
První cesta (\texttt{rust-ir}) používá přímočarý generátor
\texttt{IRRustGenerator} (\texttt{src/ir\_codegen.rs}), který převádí IR do
samostatného spustitelného programu v~Rust. Vygenerovaný kód obsahuje strukturu
\texttt{OptimizedForth} s~datovým zásobníkem \texttt{Vec<i32>} a~metodami,
které implementují jednotlivá uživatelská slova jako vnořené funkce.

Pro instrukce řídicího toku generátor používá jednoduchý stavový automat
nad čítačem instrukcí (\texttt{\_\_pc}), což umožňuje reprezentovat skoky
\texttt{Jump}/\texttt{JumpIf}/\texttt{JumpIfNot} i~bez explicitních \texttt{goto}.

Druhá cesta je novější modulární framework (\texttt{src/codegen/*}), který
odděluje překlad instrukcí (\texttt{IRTranslator}) od emitování kódu
(\texttt{CodeEmitter}) a~od specifik cílového jazyka (\texttt{TargetLanguage}).
V~této práci je tento framework využit především pro experimenty a~rozšiřitelnost.

\subsubsection{Backend pro C}

Backend pro jazyk C (\texttt{c-ir}) je implementován generátorem
\texttt{IRCGenerator} v~\texttt{src/ir\_codegen.rs}. Vygenerovaný kód
implementuje zásobník jako pole pevné velikosti \texttt{STACK\_SIZE} a~používá
pomocné funkce \texttt{push}/\texttt{pop}. Tato varianta je vhodná pro jednodušší
portaci a~snadné spuštění na systémech bez toolchainu pro Rust.

Stejně jako u~Rust backendu je možné použít i~modulární variantu backendu
(\texttt{c-modular}), která je registrována v~\texttt{codegen/registry.rs}.
